% ── §2.1: Μηχανική Μάθηση και Ανάλυση Σημάτων ───────────────

\section{Μηχανική Μάθηση και Ανάλυση Σημάτων}
\label{sec:ml_signals}

Η σύγχρονη ανάλυση σημάτων αισθητήρων βασίζεται σε μεθόδους βαθιάς μάθησης
(deep learning), οι οποίες μαθαίνουν αυτόματα ιεραρχικές αναπαραστάσεις
από τα ακατέργαστα δεδομένα, χωρίς να απαιτείται χειροκίνητη μηχανική
χαρακτηριστικών (feature engineering) \cite{Goodfellow2016}. Στο πλαίσιο
της διπλωματικής αυτής, ενδιαφερόμαστε ειδικά για την ανάλυση
\textit{πολυκαναλικών χρονοσειρών} (multivariate time-series) από
επιταχυνσιόμετρα και γυροσκόπια — δεδομένα δηλαδή με έντονη τοπική
χρονική δομή και χαμηλή εξάρτηση μεγάλης εμβέλειας.

\subsection{Συνελικτικά Νευρωνικά Δίκτυα για Χρονοσειρές}
\label{subsec:cnn_timeseries}

Τα Συνελικτικά Νευρωνικά Δίκτυα (Convolutional Neural Networks, CNN)
αρχικά αναπτύχθηκαν για την ανάλυση εικόνων, αξιοποιώντας δισδιάστατα
φίλτρα (2D kernels) για την εξαγωγή χωρικών χαρακτηριστικών \cite{He2016}.
Ωστόσο, όταν τα δεδομένα εισόδου είναι χρονοσειρές — δηλαδή ακολουθίες
μετρήσεων κατά τον χρονικό άξονα — η εφαρμογή μονοδιάστατης συνέλιξης
(1D convolution) αποδεικνύεται αποτελεσματικότερη και υπολογιστικά πιο
αποδοτική \cite{Chen2021}.

Ένα \textbf{1D-CNN επίπεδο} εφαρμόζει ένα σύνολο από $F$ φίλτρα μήκους
$k$ κατά μήκος της χρονικής διάστασης. Για είσοδο $\mathbf{X} \in
\mathbb{R}^{C \times T}$ (όπου $C$ είναι ο αριθμός καναλιών και $T$ τα
χρονικά βήματα), η έξοδος του φίλτρου $f$ στη θέση $t$ υπολογίζεται ως:

\begin{equation}
    \label{eq:conv1d}
    h_f(t) = \sigma\!\left(\sum_{c=1}^{C} \sum_{\tau=0}^{k-1}
    W_{f,c,\tau} \cdot X_{c,\, t+\tau} + b_f\right)
\end{equation}

\noindent όπου $W_{f,c,\tau}$ τα εκπαιδευόμενα βάρη του φίλτρου $f$
για κανάλι $c$ και χρονική θέση $\tau$, $b_f$ η μεροληψία (bias) και
$\sigma(\cdot)$ η συνάρτηση ενεργοποίησης ReLU: $\sigma(x) = \max(0, x)$.

Κρίσιμο πλεονέκτημα του 1D-CNN έναντι εναλλακτικών αρχιτεκτονικών,
όπως τα δίκτυα LSTM \cite{Hochreiter1997}, είναι η \textit{κοινή χρήση
παραμέτρων} (parameter sharing) κατά τον χρονικό άξονα: το ίδιο φίλτρο
εφαρμόζεται σε κάθε χρονική θέση, μειώνοντας δραστικά τον αριθμό
παραμέτρων και επιταχύνοντας την εκπαίδευση. Επιπλέον, τα 1D-CNNs
εκμεταλλεύονται φυσικά την τοπική δομή των σημάτων κίνησης (π.χ.\ μια
βηματική κύκλος διαρκεί ~0.5–1 δευτερόλεπτα, κάτι που αποτυπώνεται
σε ένα παράθυρο μερικών δεκάδων δειγμάτων).

\subsection{Δεξαμενή Χαρακτηριστικών και Πλήρης Σύνδεση}
\label{subsec:pooling_fc}

Μετά από κάθε συνελικτικό επίπεδο, εφαρμόζεται \textbf{Max Pooling}
κατά μήκος της χρονικής διάστασης, το οποίο μειώνει την ανάλυση της
χρονικής αλληλουχίας κατά παράγοντα 2 (ή μεγαλύτερο), διατηρώντας
μόνο την ισχυρότερη ενεργοποίηση σε κάθε παράθυρο. Αυτό προσφέρει
\textit{μετατοπιστική αναλλοίωτη} (translation invariance): το δίκτυο
αναγνωρίζει τοπικά μοτίβα κίνησης ανεξάρτητα από την ακριβή χρονική
θέση τους μέσα στο παράθυρο.

Στο τελικό στάδιο, ο τρισδιάστατος χάρτης χαρακτηριστικών (feature map)
ισοπεδώνεται (flatten) σε διάνυσμα και οδηγείται σε πλήρως συνδεδεμένα
(fully-connected) επίπεδα, τα οποία παράγουν λογιτικές βαθμολογίες
ανά κλάση δραστηριότητας. Η τελική ταξινόμηση γίνεται μέσω \textit{softmax}:

\begin{equation}
    \label{eq:softmax}
    p_j = \frac{e^{z_j}}{\sum_{l=1}^{K} e^{z_l}}, \quad j = 1,\ldots,K
\end{equation}

\noindent όπου $z_j$ οι λογιτικές βαθμολογίες και $K = 6$ ο αριθμός
κλάσεων δραστηριότητας.

\subsection{Σύγκριση με Εναλλακτικές Αρχιτεκτονικές}
\label{subsec:architecture_comparison}

Η ολοκληρωμένη ανασκόπηση των Chen et al.\ \cite{Chen2021} αξιολογεί
συστηματικά CNN, LSTM και υβριδικές αρχιτεκτονικές για HAR. Τα
ευρήματά τους δείχνουν ότι τα 1D-CNN επιτυγχάνουν ανταγωνιστική ακρίβεια
συγκριτικά με τα LSTM, ενώ παρουσιάζουν σημαντικά μικρότερο υπολογιστικό
κόστος — ιδιότητα κρίσιμη σε FL περιβάλλοντα όπου η τοπική εκπαίδευση
γίνεται σε περιορισμένους πόρους (smartphones). Πιο πρόσφατα, μηχανισμοί
\textit{self-attention} (Transformers) \cite{Vaswani2017} έχουν εφαρμοστεί
και σε HAR χρονοσειρές, επιτυγχάνοντας υψηλή ακρίβεια, ωστόσο ο
σημαντικά μεγαλύτερος αριθμός παραμέτρων τους τα καθιστά λιγότερο
κατάλληλα για τοπική εκπαίδευση σε περιβάλλον FL με περιορισμένους πόρους.
Για τον λόγο αυτό, η παρούσα εργασία υιοθετεί αρχιτεκτονική 1D-CNN, η
οποία περιγράφεται αναλυτικά στην Ενότητα~\ref{sec:cnn_design}.
